\chapter{固体物理基础}
\section{从波函数到宏观物理量}
凝聚态物理的起点是多粒子的薛定谔方程,$\Psi(\{ \mbf{r}_i\},t)$是反对称的多粒子波函数，简单起见我们暂时不考虑自旋。
\begin{align}
    i\hbar\dv{\Psi(\{ \mathbf{r}_i\},t)}{t}=\hat{H}\Psi(\{ \mathbf{r}_i\},t)
    \label{hamiltonian}
\end{align}
哈密顿算符是位置算符、动量算符（以及时间和系统所处外部环境）的函数，其函数形式与经典力学的哈密顿量相同，包含了系统与相互作用相关的全部信息。
\begin{align}
    \hat{H}=-\frac{\hbar^2}{2m_e}\sum_i \nabla_{i}^{2}-\frac{\hbar^2}{2M_I}\sum_I \nabla_{I}^{2}-\sum_{i,I}\frac{Z_Ie^2}{|r_i-R_I|}+\frac{1}{2}\sum_{i\neq j}\frac{e^2}{r_i-r_j}+\frac{1}{2}\sum_{I\neq J}\frac{Z_IZ_Je^2}{|R_I-R_J|}
\end{align}
其中$Z_I$和$M_I$是原子核的电荷数和质量。
\subsection{电子密度}
电子密度算符:
\begin{gather}
    \hat{n}(\mathbf(r))=\sum_{i=1}^{N}\delta(\mathbf{r}-\mathbf{r}_i)
    \end{gather}
    期望值:
    \begin{gather}
    n(\mathbf{r})=\frac{\langle\Psi|\hat{n}(\mathbf{r}\rangle)}{\langle\Psi|\Psi\langle}
    =N\frac{\int \dd^3r_2\dd^3r_3\cdots\dd^3r_N|\Psi({r,r_2,r_3\cdots r_N})|^2}{\int \dd^3r_1\dd^3r_2\dd^3r_3\cdots\dd^3r_N|\Psi({r_1,r_2,r_3\cdots r_N})|^2}
\end{gather}
\subsection{能量}
忽略掉原子核的动能，系统的哈密顿量可以分为以下几部分：
\begin{equation}
    \hat{H}=\hat{T}+\hat{V}_{\text{ext}}+\hat{V}_{\text{int}}+E_{II}
\end{equation}
它们分别对应于公式电子动能、电子与核的相互作用、电子之间的相互作用、原子核之间的相互作用。$\hat{H}$的期望值也可以拆成这几项：
\begin{equation}
    E\equiv \langle\hat{H}\rangle=\frac{\Bk{\hat{H}}}{\langle\Psi|\Psi\rangle}=\langle\hat{T}\rangle+\langle\hat{V}_{\text{int}}\rangle+\int \dd^3r V_{\text{ext}}(\mathbf{r})n(\mathbf{r})+E_{II}
    \label{energy}
\end{equation}
如果将能量E看做是波函数$\Psi$的泛函，那么哈密顿量的本征态就是使能量取得极值的鞍点。我们用拉格朗日乘子法处理公式\ref{energy}，得到的结果是定态薛定谔方程：
\begin{equation}
    \hat{H}|\Psi\rangle=E|\Psi\rangle
\end{equation}
\section{晶体结构}
\begin{itemize}
    \item 理想的晶体是由相同的\tbf{原子团}（注意是原子团而不是原子）在空间中\tbf{无限重复}排列构成的，这样的原子团称为\tbf{基元（basis）}。
    \item 在研究晶体的结构时，这些原子团可以抽象为几何点，所有点的集合称为\tbf{晶格（lattice）}。
    \item 由于晶体具有平移对称性，可以通过平移$u_1\vec{a}_1+u_2\vec{a}_2+u_3\vec{a}_3$使每一个格点移动到完全等价的另一个格点。如果以上的操作可以遍历全部格点，那么这个晶格就称为初基晶格，$\vec{a}_i$称为\tbf{初基平移矢量(primitive translation vector)}。基元内原子j的中心位置也可以用$\vec{a}$来表示：
    \begin{equation*}
        \mathbf{r}_j=x_j\mathbf{a}_1+y_j\mathbf{a}_2+z_j\mathbf{a}_3
    \end{equation*}初基平移矢量的选取不是唯一的。我们往往用初基平移矢量来定义晶轴。
    \item 体积最小的晶胞称为原胞。初基平移矢量所确定的平行六面体($\vec{a}_1\cdot\vec{a}_2\times\vec{a}_3$)是原胞的一种选取方法，因为初级平移矢量的选取不唯一，原胞选取也不唯一。
    \item \tbf{维格纳-赛茨原胞}：原胞的另一种取法，把一个格点和相邻格点用直线连接起来，做连线的垂直平分线或者垂直平分面，围成的最小体积。
\end{itemize}
\section{倒空间}
周期性的函数如晶体中的势能、电子密度等可以展开成傅里叶级数，以电子密度为例:
\begin{align}
    n(\mathbf{r})&=\sum_{\mathbf{G}}n_Ge^{i\mathbf{G}\cdot\mathbf{r}}\\
    n_G&=\frac{1}{V_c}\int_{cell}n(\vec{r})e^{-i\mathbf{G\cdot r}}\dd{\vec{r}}
\end{align}
$\mathbf{G}$不是任意的,因为$n(\mathbf{r})$具有平移不变性:
\begin{equation}
    n(\mathbf{r}+\vec{T})=\sum_{\mathbf{G}}n_Ge^{i\mathbf{G}\cdot\mathbf{r}}e^{i\mathbf{G}\cdot\mathbf{T}}=n(\mathbf{r})
\end{equation}
我们定义倒格子的轴矢量:
\begin{equation}
    \mathbf{b}_1=2\pi\frac{\mathbf{a}_2\times\mathbf{a}_3}{\mathbf{a_1}\cdot\mathbf{a}_2\times\mathbf{a}_3}
    \quad
    \mathbf{b}_2=2\pi\frac{\mathbf{a}_3\times\mathbf{a}_1}{\mathbf{a_2}\cdot\mathbf{a}_3\times\mathbf{a}_1}
    \quad
    \mathbf{b}_3=2\pi\frac{\mathbf{a}_1\times\mathbf{a}_2}{\mathbf{a_3}\cdot\mathbf{a}_1\times\mathbf{a}_3}
\end{equation}
倒空间每个初基矢量和实空间的两个轴矢量正交:
\begin{equation}
    \mathbf{b}_i\cdot\mathbf{a}_j=2\pi\delta_{ij}
\end{equation}
当$\mathbf{G}=v_1\mathbf{b}_1+v_2\mathbf{b}_2+v_3\mathbf{b}_3(v_i$是整数)时,$\exp(i\mathbf{G}\cdot\mathbf{T})=1$,因此$\mathbf{G}=v_1\mathbf{b}_1+v_2\mathbf{b}_2+v_3\mathbf{b}_3$就是波矢被允许的取值范围.
\subsection{布里渊区}
布里渊区的定义是倒格子空间中的维格纳-赛茨原胞,它的边界是服从衍射条件的波矢$\mathbf{k}$:
\begin{equation}
    2\mathbf{k}\cdot\mathbf{G}=G^2
\end{equation}
倒格子的中央晶胞称为第一布里渊区.
\section{自由电子气}
\section{固体能带理论}
\subsection{波恩-冯卡门边界条件}
\subsection{近自由电子近似}
近自由电子近似将晶体中离子实对电子的作用,看过是一个周期势场微扰.一个具有平移操作不变性的周期性势场可以展开成倒格矢的傅里叶级数:
\begin{equation}
    U(x)=\sum_G U_Ge^{iGx}
\end{equation}
\subsection{布洛赫定理}
平移算符和哈密顿算符对易,二者具有相同的本征函数:
\begin{align}
    \hat{H}\psi(r)=E\psi(r) \qquad \hat{T}(r)\psi(r)=C(r)\psi(r)
\end{align}
平移操作满足交换律$\hat{T}(r_1)\hat{T}(r_2)=\hat{T}(r_2)\hat{T}(r_1)=\hat{T}(r_1+r_2)$,因此布拉维平移群是阿贝尔群,根据舒尔引理,平移算符的本征值 C(r) 构成一维表示.
由于群乘法规则$\hat{T}(r_1)\hat{T}(r_2)=\hat{T}(r_1+r_2)$,本征值需要满足:$C(r_1)C(r_2)=C(r_1+r_2)$,唯一连续解是指数函数:
\begin{equation}
    C(r)=e^{ik\cdot r}
\end{equation}
于是$\hat{T}(a)\psi(r)=\psi(r+a)=e^{ik\cdot a}\psi(r)$(注意不是级数).
为了满足这一关系，波函数必须包含一个随位置变化的相位因子$e^{ik\cdot r}$,以抵消平移带来的相位变化.将这部分显式地提取出来:
\begin{gather}
    \psi(r)=e^{ik\cdot r}u(r)\rightarrow
    \psi(r+a)=e^{ik\cdot(r+a)}u(r+a)=e^{ik\cdot a}e^{ik\cdot r}u(r)
\end{gather}
因此u(r)是周期为a的函数:
\begin{gather}
    u(r+a)=u(r)
\end{gather}
\subsection{电子在周期性势场的波动方程}
考虑一个晶格常数为a,首尾相连(波恩-冯卡门边界条件),周长为L的一维晶格圆环.单电子近似下,晶体中电子的波动方程可以写为以下形式:
\begin{equation}
    (\frac{1}{2m}p^2+U(x))\psi(x)=\epsilon\psi(x)
    \label{seq_c}
\end{equation}
我们将单电子波函数$\psi(x)$展开,
\begin{equation}
    \psi(x)=\sum_{k}C_{k}e^{ikx}
\end{equation}
波矢k需要满足$k=\frac{2\pi n}{L}\text{(n是整数)}$,因为$\psi(x)$是周期为L的周期性函数(参阅傅里叶级数).布洛赫定理对k
提出了额外的要求:
\[
    \left\{\begin{array}{l}
        \psi(x)=e^{ikx}\sum_{k'}C_ke^{i(k'-k)x}\\
        \psi(x+a)=e^{ik(x+a)}\sum_{k'}C_ke^{i(k'-k)x}e^{i(k'-k)a}
    \end{array}\right .\rightarrow e^{i(k'-k)a}=1\rightarrow k'=k+G
    \]
因此一旦波矢k出现在$\psi(x)$的展开式中,那么其它展开项只能是符合k+G的波矢.将势能和波函数的级数代入
\label{seq_c},得到:
\begin{align}
\sum_{k}\frac{\hbar^2k^2}{2m}C_ke^{ikx}+\sum_{G}\sum_{k}U_GC_ke^{i(k+G)x}=\epsilon\sum_{k}C_ke^{ikx}
\end{align}
方程两端各个k的傅里叶分量系数相等:
\begin{gather}
    (\lambda_k-\epsilon)C_k+\sum_{G}U_GC_{k-G}=0 \quad\quad\text{其中}\quad\lambda_k=\frac{\hbar^2k^2}{2m}
\end{gather}
上式称为\tbf{中心方程}.
\section{附录:三维空间的十四种晶格}
\section{附录:晶体衍射条件}